\documentclass{report}
\usepackage[utf8]{inputenc}
\usepackage[T1]{fontenc}
\usepackage[french]{babel}
\usepackage{geometry}
\usepackage{fancyhdr}
\setlength{\headheight}{13pt}
\usepackage{lastpage}
\usepackage{titlesec}
\usepackage{listings}
\usepackage{enumitem}
\usepackage{color}
\usepackage{xcolor}
\usepackage{soul}
\usepackage{graphicx}
\usepackage[utf8]{inputenc}
\usepackage{amsmath}
\usepackage{hyperref}
\hypersetup{hidelinks}
\usepackage{tikz}
\usetikzlibrary{positioning}
\usepackage{longtable}



\newcommand{\DocName}{OPTIMISATION DE MODELES NEURONAUX A OSCILLATEURS COUPLÉS - LIVRABLE 1}
\newcommand{\DocAuthors}{\bsc{Da Silva} Tristan - \bsc{Battaglia} Charles-Antoine \\ \bsc{Botté-Magalhaes} Jules - \bsc{Gauthier} Keanu - \bsc{Fraysse} Ianis}

\definecolor{dkgreen}{rgb}{0,0.6,0}
\definecolor{gray}{rgb}{0.5,0.5,0.5}
\definecolor{mauve}{rgb}{0.58,0,0.82}
\definecolor{addrred}{rgb}{0.93,0.43,0.43}
\definecolor{commentgray}{rgb}{0.45,0.45,0.45}

\lstdefinelanguage{MicroAsm}{
    sensitive=true,
    morecomment=[l]{\#},
    moredelim=[s][\color{addrred}]{@}{:}
}

\lstdefinelanguage{HexProg}{
    sensitive=true,
    morecomment=[l]{;}
}

\lstdefinelanguage{AsmStd}{
    sensitive=true,
    morecomment=[l]{;},
    morecomment=[l]{\#},
    morekeywords={
        MOV,ADD,SUB,MUL,DIV,INC,DEC,CMP,AND,OR,XOR,NOT,
        JMP,JE,JNE,JG,JL,JGE,JLE,CALL,RET,PUSH,POP,
        LDA,STA,LOAD,STORE,INPUT,OUTPUT,HALT,NOP
    }
}

\lstset{frame=tb,
    language=MicroAsm,
    aboveskip=3mm,
    belowskip=3mm,
    showstringspaces=false,
    columns=flexible,
    basicstyle={\small\ttfamily},
    numbers=none,
    numberstyle=\tiny\color{gray},
    keywordstyle=\color{blue},
    commentstyle=\color{commentgray}, 
    breaklines=true,
    breakatwhitespace=true,
    tabsize=3,
   }

\lstset{breaklines=true}

\geometry{a4paper, margin=1in}

% Suppression de l'affichage du mot "Chapitre"
\titleformat{\chapter}[display]
    {\normalfont\bfseries\LARGE}
    {}
    {0pt}
    {}

\definecolor{Light}{gray}{.90}
\sethlcolor{Light}

\let\OldTexttt\texttt
\renewcommand{\texttt}[1]{\OldTexttt{\hl{#1}}}

\renewcommand{\thechapter}{\arabic{chapter}}
\renewcommand{\thesection}{\thechapter.\arabic{section}}
\renewcommand{\thesubsection}{\thesection.\arabic{subsection}}

\fancyhf{}
\fancyhead[L]{\DocName}
\fancyfoot[L]{\DocAuthors}
\fancyfoot[R]{\thepage / \pageref{LastPage}}
\renewcommand{\headrulewidth}{0 pt}
\renewcommand{\footrulewidth}{0.2pt}
\pagestyle{fancy}
\fancypagestyle{plain}{
    \fancyhf{}
    \fancyfoot[R]{\thepage / \pageref{LastPage}}
    \fancyfoot[L]{\DocAuthors}
    \fancyhead[L]{\DocName}
}

\graphicspath{ {./img} }

\title{\DocName}
\author{\DocAuthors}
\date{\textbf{JUNIA - ISEN - AP4} \url{https://github.com/KeanuGauthier/Socle-Modele-Neuronaux-RD.git}}

\newcounter{questioncounter}
\setcounter{questioncounter}{0}

\begin{document}
\maketitle
\tableofcontents

\chapter{Contexte et motivation}

Ce projet s'inscrit dans le cadre du laboratoire de recherche de M.~Antoine Pirog, encadrant du projet R\&D au semestre 7 à Junia. Le laboratoire travaille sur la simulation de réseaux de neurones biophysiques, c'est-à-dire des modèles qui reproduisent le comportement électrique de populations de neurones interconnectés.

Le problème auquel ce type de recherche fait face est simple : la simulation de réseaux de neurones réalistes est coûteuse en temps de calcul. Sur CPU, un réseau de quelques milliers de neurones peut prendre des heures à simuler. Cela limite la taille des réseaux qu'on peut étudier, et donc les questions scientifiques qu'on peut explorer.

Junia dispose de serveurs équipés de cartes graphiques RTX~5090. La question posée est donc : ces GPU peuvent-ils apporter la puissance de calcul nécessaire pour rendre la simulation de grands réseaux de neurones tractable au quotidien pour le laboratoire ?

La réponse n'est pas évidente, et elle n'est pas garantie. Le GPU est théoriquement beaucoup plus puissant, mais notre problème, un système d'EDO couplées avec des connexions irrégulières, n'est pas forcément un cas idéal pour le parallélisme massif. Il est tout à fait possible que la conclusion soit : non, sur ce type de simulation, le GPU n'apporte pas un gain suffisant pour justifier le changement d'architecture. Le but du projet est d'apporter une réponse documentée et mesurée à cette question, pas d'arriver à une conclusion prédéterminée.

C'est dans ce contexte que le sujet \og Optimisation de modèles neuronaux à oscillateurs couplés \fg{} nous a été confié, avec un socle de code déjà existant : modèle de Hodgkin-Huxley, synapses, solveur CPU de référence. On part de cette base pour évaluer l'intérêt du GPU.

\chapter{Présentation du sujet}

\section{Énoncé et objectif}

Ce projet a pour objectif d'évaluer la faisabilité de l'utilisation des architectures GPU pour la simulation de réseaux de neurones. Il s'agit de déterminer la capacité maximale du système en nombre de synapses prises en charge, tout en précisant les conditions dans lesquelles le calcul sur GPU surpasse une exécution sur CPU. Enfin, les travaux visent à accélérer et à optimiser les performances de simulation en s'appuyant sur nos compétences techniques.

L'approche suivie s'articule en cinq étapes : modélisation Hodgkin-Huxley, couplage par synapses, construction d'un connectome, portage sur GPU, puis benchmark CPU/GPU. L'objectif final est de produire une courbe de speedup en fonction de la taille du réseau, qui montre à partir de quel nombre de neurones le GPU devient indispensable, et dans quelle mesure il l'est.

\section{Le modèle de Hodgkin-Huxley}

Le socle de code fourni repose sur le modèle de Hodgkin-Huxley. Il s'agit d'une modélisation biologique dans laquelle un neurone est assimilé à un circuit électrique, simulant l'activité cérébrale à partir du flux d'ions sodium et potassium. Le modèle décrit l'évolution du potentiel membranaire par un système d'équations différentielles ordinaire (EDO) gouvernant l'ouverture et la fermeture des canaux ioniques. C'est ce formalisme, une EDO par neurone répétée sur l'ensemble du réseau, qui rend le problème naturellement parallélisable.

\section{Problématique}

\textbf{Problématique.} Comment accélérer la simulation de réseaux de neurones de Hodgkin-Huxley couplés lorsque la taille du connectome augmente, et quelle stratégie GPU offre le meilleur gain par rapport à la référence CPU, à résultats équivalents ?

\textbf{La vraie question du projet.} Le laboratoire de M.~Pirog peut-il, grâce aux serveurs Junia équipés de RTX~5090, simuler des réseaux de neurones assez grands pour être utiles à la recherche, dans un temps acceptable ?

\end{document}